%======================================================================
% Chapter 4: Linear Algebra I: Vectors and Matrices
%======================================================================
\chapter{Linear Algebra I: Vectors and Matrices}\label{ch:linalg1}

\begin{motivation}
Open the source code of any machine learning system and you will find the
same two kinds of object everywhere: \emph{vectors} (lists of numbers) and
\emph{matrices} (tables of numbers). A data point is a vector of features; a
data set is a matrix with one row per data point; an image is a matrix of
pixels; the weights of a neural network layer form a matrix; and the
computation done by that layer is ``multiply the input vector by the weight
matrix, add a bias vector, apply an activation function''. Graphics processing
units (GPUs), the hardware on which modern AI runs, are essentially machines
for multiplying matrices very fast.

Linear algebra, the mathematics of vectors and matrices, therefore gives us
two things at once: a compact \emph{notation} in which a single short formula
describes millions of arithmetic operations, and a geometric \emph{picture} in
which a matrix is a transformation that stretches, rotates and squashes space.
This chapter covers the basics; \cref{ch:linalg2} covers the deeper ideas that
machine learning relies on.
\end{motivation}

\begin{objectives}
\begin{itemize}
  \item represent data as vectors, and picture vectors as arrows;
  \item add and scale vectors and form linear combinations;
  \item compute lengths, distances, dot products and angles between vectors;
  \item store data sets and images as matrices;
  \item multiply a matrix by a vector and by another matrix, and explain both
    operations in two complementary ways;
  \item interpret a matrix as a linear transformation of space;
  \item use the transpose, the identity matrix, inverses and $2\times2$ and
    $3 \times 3$ determinants;
  \item write a system of linear equations as $\mA\vx = \vb$ and solve it by
    Gaussian elimination.
\end{itemize}
\end{objectives}

%----------------------------------------------------------------------
\section{Vectors}\label{sec:vectors}
%----------------------------------------------------------------------

\subsection{Two ways to see a vector}

The word \emph{vector} is used for two ideas that turn out to be the same
thing.

\begin{itemize}
  \item \textbf{A vector is an arrow.} In physics a vector is a quantity with a
    size and a direction, such as a displacement (``$3$\,m east and $2$\,m
    north'') or a velocity. We draw it as an arrow. Only its length and
    direction matter, not where it is drawn: the three arrows in
    \cref{fig:vector-arrow} all represent the same vector.
  \item \textbf{A vector is a list of numbers.} In computer science and data
    science a vector is an ordered list of numbers, such as the features of a
    house: $(120, 3, 15)$ for $120$\,m$^2$, $3$ bedrooms and $15$ years of age.
\end{itemize}

The bridge between the two views is the coordinate system of
\cref{ch:geometry}: the arrow that starts at the origin and ends at the point
$(3, 2)$ corresponds to the list $(3, 2)$. The numbers say how far the arrow
goes in each coordinate direction.

\begin{figure}[htbp]
\centering
\begin{tikzpicture}[scale=0.75]
  \draw[gray!22] (-0.5,-0.5) grid (7.5,5.5);
  \draw[-{Stealth[length=2.5mm]},thick] (-0.5,0) -- (7.8,0) node[right] {$x_1$};
  \draw[-{Stealth[length=2.5mm]},thick] (0,-0.5) -- (0,5.8) node[above] {$x_2$};
  \draw[dashed,gray] (3,0) -- (3,2) -- (0,2);
  \draw[vec,PlotBlue] (0,0) -- (3,2);
  \node[PlotBlue,above left] at (1.5,1) {$\vv$};
  \fill[PlotBlue] (3,2) circle (2.5pt);
  \node[right,font=\small] at (3.05,1.8) {$(3,2)$};
  \draw[vec,PlotBlue!55] (4,0.5) -- (7,2.5);
  \draw[vec,PlotBlue!55] (2,3) -- (5,5);
  \node[font=\small,below] at (1.5,0) {$3$};
  \node[font=\small,left] at (0,1) {$2$};
\end{tikzpicture}
\caption{The vector $\vv = (3, 2)$ drawn as an arrow from the origin. The two
  paler arrows have the same length and direction, so they represent the same
  vector: ``$3$ to the right and $2$ up''.}
\label{fig:vector-arrow}
\end{figure}

\begin{definition}[Vector]\label{def:vector}
A \term{vector} in $\R^n$ is an ordered list of $n$ real numbers, called its
\term[component of a vector]{components} (or entries). We write vectors as
columns:
\[
  \vv = \begin{pmatrix} v_1 \\ v_2 \\ \vdots \\ v_n \end{pmatrix}.
\]
The number $n$ is the \emph{dimension} of the space $\R^n$. Two vectors are
equal when all their components are equal.
\end{definition}

\begin{notation}
Vectors are written in \textbf{bold} lower-case letters: $\vv$, $\vx$, $\vw$.
Their components are written in ordinary italic with a subscript: $v_1$,
$x_3$, $w_i$. Single numbers are called \term[scalar]{scalars}, to distinguish
them from vectors, and are written in ordinary italic: $a$, $c$, $\lambda$.
(On paper, where bold is hard to write, people often write $\vec v$ or
underline the letter.) To save space we sometimes write a column vector on one
line as $\vv = (v_1, \dots, v_n)$; it still means a column. The
\term{zero vector} $\vzero = (0, \dots, 0)$ has every component equal to $0$.
\end{notation}

\begin{mlnote}[Everything is a vector]
In machine learning, almost every object is stored as a vector.
\begin{itemize}
  \item \textbf{Tabular data:} a house is $\vx = (120, 3, 15, 2.5)$: size,
    bedrooms, age and distance to the centre.
  \item \textbf{Images:} a $28\times28$ grayscale image becomes a vector in
    $\R^{784}$ by reading its pixels row by row (``flattening'').
  \item \textbf{Words:} language models represent each word (or word piece) by
    a learned vector called an \term{embedding}, typically with hundreds or
    thousands of components.
  \item \textbf{Categories:} a category out of $k$ possibilities is encoded as a
    \term{one-hot vector} with a single $1$: with the categories cat, dog,
    bird, ``dog'' is $(0, 1, 0)$.
  \item \textbf{Model parameters:} all the weights of a model can be listed in
    one long vector, usually called $\vtheta$ (``theta'').
\end{itemize}
\end{mlnote}

\sectionexercises

\begin{exercise}
Write a feature vector for a house of $120$\,m$^2$ with $3$ bedrooms, aged
$15$ years and $2.5$\,km from the city centre. In which space $\R^n$ does it
live?
\end{exercise}
\begin{solution}
$\vx = (120, 3, 15, 2.5)$, written as a column; it lives in $\R^4$.
\end{solution}

\begin{exercise}
Write the one-hot vectors for the categories red, green and blue (in that
order). What is the sum of the entries of any one-hot vector?
\end{exercise}
\begin{solution}
Red $= (1,0,0)$, green $= (0,1,0)$, blue $= (0,0,1)$. The entries of a one-hot
vector always add up to $1$.
\end{solution}

\begin{exercise}
Which vector is represented by the arrow from $P(1, 4)$ to $Q(5, 1)$? How long
is the arrow?
\end{exercise}
\begin{solution}
The displacement is $Q - P = (5 - 1, 1 - 4) = (4, -3)$. Its length is
$\sqrt{16 + 9} = 5$.
\end{solution}

%----------------------------------------------------------------------
\section{Adding and scaling vectors}\label{sec:vector-operations}
%----------------------------------------------------------------------

\subsection{Vector addition}

Vectors of the same dimension are added component by component:
\[
  \vu + \vv = \begin{pmatrix} u_1 \\ \vdots \\ u_n \end{pmatrix}
  + \begin{pmatrix} v_1 \\ \vdots \\ v_n \end{pmatrix}
  = \begin{pmatrix} u_1 + v_1 \\ \vdots \\ u_n + v_n \end{pmatrix}.
\]
For example, $(3, 1) + (1, 2) = (4, 3)$. Geometrically, addition means
``follow one arrow, then the other'': place the tail of $\vv$ at the tip of
$\vu$; the sum goes from the tail of $\vu$ to the tip of $\vv$. Doing it in
the other order gives the same result, and the two routes form a
parallelogram (\cref{fig:vector-ops}, left). Vectors of different dimensions
cannot be added.

\subsection{Scalar multiplication}

Multiplying a vector by a scalar multiplies every component:
\[
  c\,\vv = \begin{pmatrix} c\,v_1 \\ \vdots \\ c\,v_n \end{pmatrix}.
\]
Geometrically, $c\vv$ points along the same line as $\vv$ and is $\abs{c}$ times
as long; if $c < 0$ it points the opposite way (\cref{fig:vector-ops},
right). The vector $(-1)\vv$ is written $-\vv$, and subtraction is defined by
$\vu - \vv = \vu + (-\vv)$, i.e.\ component by component. Geometrically,
$\vu - \vv$ is the arrow from the tip of $\vv$ to the tip of $\vu$.

\begin{figure}[htbp]
\centering
\begin{subfigure}{0.48\textwidth}
\centering
\begin{tikzpicture}[scale=0.85]
  \draw[gray!22] (-0.4,-0.4) grid (4.6,3.6);
  \draw[-{Stealth[length=2mm]}] (-0.4,0) -- (4.8,0);
  \draw[-{Stealth[length=2mm]}] (0,-0.4) -- (0,3.8);
  \draw[vec,PlotBlue] (0,0) -- (3,1) node[midway,below right] {$\vu$};
  \draw[vec,PlotRed] (0,0) -- (1,2) node[midway,above left] {$\vv$};
  \draw[vec,PlotRed!55,dashed] (3,1) -- (4,3) node[midway,right] {$\vv$};
  \draw[vec,PlotBlue!55,dashed] (1,2) -- (4,3) node[midway,above] {$\vu$};
  \draw[vec,PlotGreen] (0,0) -- (4,3) node[pos=0.62,above left] {$\vu+\vv$};
\end{tikzpicture}
\caption{The parallelogram rule: $(3,1) + (1,2) = (4,3)$.}
\end{subfigure}\hfill
\begin{subfigure}{0.48\textwidth}
\centering
\begin{tikzpicture}[scale=0.75]
  \draw[gray!22] (-2.4,-1.4) grid (4.4,2.4);
  \draw[-{Stealth[length=2mm]}] (-2.4,0) -- (4.6,0);
  \draw[-{Stealth[length=2mm]}] (0,-1.4) -- (0,2.6);
  \draw[vec,PlotGreen!80!black] (0,0) -- (4,2) node[above] {$2\vv$};
  \draw[vec,PlotBlue] (0,0) -- (2,1) node[below right] {$\vv$};
  \draw[vec,PlotOrange] (0,0) -- (1,0.5) node[above left] {$\tfrac12\vv$};
  \draw[vec,PlotRed] (0,0) -- (-2,-1) node[left] {$-\vv$};
\end{tikzpicture}
\caption{Scalar multiples of $\vv = (2,1)$.}
\end{subfigure}
\caption{The two basic vector operations.}
\label{fig:vector-ops}
\end{figure}

These operations obey the same laws as ordinary arithmetic
(\cref{thm:laws-arithmetic}): $\vu + \vv = \vv + \vu$,
$(\vu + \vv) + \vw = \vu + (\vv + \vw)$, $\vv + \vzero = \vv$,
$\vv + (-\vv) = \vzero$, $c(\vu + \vv) = c\vu + c\vv$, $(c + d)\vv = c\vv +
d\vv$ and $c(d\vv) = (cd)\vv$. Any collection of objects with an addition and a
scalar multiplication obeying these laws is called a \term{vector space}; $\R^n$
is the main example.

\subsection{Linear combinations}

Combining the two operations gives the single most important construction in
linear algebra.

\begin{definition}[Linear combination]\label{def:linear-combination}
A \term{linear combination} of vectors $\vv_1, \dots, \vv_k$ is any vector of
the form
\[
  c_1\vv_1 + c_2\vv_2 + \dots + c_k\vv_k ,
\]
where $c_1, \dots, c_k$ are scalars, called the \emph{coefficients} or
\emph{weights} of the combination.
\end{definition}

\begin{example}[A linear combination]\label{ex:linear-combination}
Compute $2\vu - 3\vv$ for $\vu = (1, 0, 2)$ and $\vv = (2, -1, 1)$.

\sol $2\vu = (2, 0, 4)$ and $3\vv = (6, -3, 3)$, so
$2\vu - 3\vv = (2 - 6, 0 + 3, 4 - 3) = (-4, 3, 1)$.
\end{example}

The vectors $\ve_1 = (1, 0)$ and $\ve_2 = (0, 1)$ are special: every vector in
$\R^2$ is a linear combination of them, in exactly one way:
\[
  \begin{pmatrix} 3 \\ 2 \end{pmatrix}
  = 3\begin{pmatrix} 1 \\ 0 \end{pmatrix}
  + 2\begin{pmatrix} 0 \\ 1 \end{pmatrix}
  = 3\ve_1 + 2\ve_2 .
\]
In $\R^n$, the \term[standard basis]{standard basis vectors}
$\ve_1, \dots, \ve_n$ have a $1$ in one position and $0$ elsewhere, and every
vector is $\vv = v_1\ve_1 + \dots + v_n\ve_n$. The components of a vector are
just the coefficients needed to build it from the standard basis.

\begin{intuition}[Linear combinations are recipes]
Think of $\vv_1, \dots, \vv_k$ as ingredients and of the coefficients as
amounts. If one portion of food A provides $(10, 5, 30)$ grams of (protein,
fat, carbohydrate) and one portion of food B provides $(20, 10, 5)$, then two
portions of A and one of B provide $2(10, 5, 30) + (20, 10, 5) = (40, 20,
65)$. Which nutrient totals can you reach by mixing A and B? That question,
``which vectors are linear combinations of these?'', is the subject of
\cref{ch:linalg2}.
\end{intuition}

\begin{mlnote}[Vector arithmetic in machine learning]
Adding and scaling vectors is the bread and butter of machine learning code.
\begin{itemize}
  \item The \emph{mean} of data vectors $\vx_1, \dots, \vx_n$ is the linear
    combination $\frac1n(\vx_1 + \dots + \vx_n)$, computed component by
    component. It is the ``centre of mass'' of the data cloud.
  \item Training updates all weights at once with a vector operation: gradient
    descent replaces $\vtheta$ by $\vtheta - \eta\,\vg$, where $\vg$ is a vector
    (the gradient, \cref{ch:multivariable}) and $\eta$ a small scalar
    (\cref{ch:optimization}).
  \item Word embeddings famously support analogies by vector arithmetic: the
    vector $\text{king} - \text{man} + \text{woman}$ lies close to
    $\text{queen}$.
\end{itemize}
\end{mlnote}

\sectionexercises

\begin{exercise}
Let $\vu = (2, -1, 3)$ and $\vv = (-1, 4, 0)$. Compute $\vu + \vv$,
$\vu - \vv$, $3\vu$ and $2\vu - 3\vv$.
\end{exercise}
\begin{solution}
$\vu + \vv = (1, 3, 3)$; $\vu - \vv = (3, -5, 3)$; $3\vu = (6, -3, 9)$;
$2\vu - 3\vv = (4, -2, 6) - (-3, 12, 0) = (7, -14, 6)$.
\end{solution}

\begin{exercise}
Find scalars $a$ and $b$ such that $a(1, 1) + b(1, -1) = (5, 1)$.
\end{exercise}
\begin{solution}
Component by component: $a + b = 5$ and $a - b = 1$. Adding gives $2a = 6$, so
$a = 3$ and $b = 2$. Check: $3(1,1) + 2(1,-1) = (5, 1)$.
\end{solution}

\begin{exercise}
Compute the mean of the data vectors $(1, 2)$, $(3, 6)$ and $(5, 1)$.
\end{exercise}
\begin{solution}
$\tfrac13\bigl((1,2) + (3,6) + (5,1)\bigr) = \tfrac13(9, 9) = (3, 3)$.
\end{solution}

\begin{exercise}
A model has weights $\vtheta = (0.5, -1.2, 2.0)$. One step of gradient descent
with learning rate $\eta = 0.1$ and gradient $\vg = (0.4, -0.6, 1.0)$ replaces
$\vtheta$ by $\vtheta - \eta\vg$. Compute the new weights.
\end{exercise}
\begin{solution}
$\eta\vg = (0.04, -0.06, 0.1)$, so the new weights are
$(0.5 - 0.04, -1.2 + 0.06, 2.0 - 0.1) = (0.46, -1.14, 1.9)$.
\end{solution}

%----------------------------------------------------------------------
\section{Length, distance and unit vectors}\label{sec:vector-length}
%----------------------------------------------------------------------

The length of the arrow $\vv = (3, 4)$ is the hypotenuse of a right triangle
with legs $3$ and $4$, namely $5$. In general, by the Pythagorean theorem in
$n$ dimensions (\cref{def:euclidean-distance}):

\begin{definition}[Length (norm) of a vector]\label{def:norm}
The \term{length}, or \term[norm!Euclidean]{Euclidean norm}, of
$\vv \in \R^n$ is
\[
  \norm{\vv} = \sqrt{v_1^2 + v_2^2 + \dots + v_n^2}.
\]
The \emph{distance} between $\vu$ and $\vv$ is the length of their
difference, $\norm{\vu - \vv}$.
\end{definition}

For example, $\norm{(1, -2, 2)} = \sqrt{1 + 4 + 4} = 3$. Lengths behave as
expected under scaling: $\norm{c\vv} = \abs{c}\,\norm{\vv}$. (Doubling a vector
doubles its length; reversing it does not change its length.)

A vector of length $1$ is a \term{unit vector}. Any non-zero vector can be
turned into a unit vector pointing in the same direction by dividing by its
length, which is called \term[normalisation of a vector]{normalising} it:
\[
  \hat\vv = \frac{\vv}{\norm{\vv}} .
\]
For $\vv = (3, 4)$, $\hat\vv = (0.6, 0.8)$. A unit vector records a pure
direction.

\begin{theorem}[Triangle inequality]\label{thm:triangle-inequality}
For all vectors $\vu, \vv \in \R^n$,
\[
  \norm{\vu + \vv} \le \norm{\vu} + \norm{\vv} .
\]
\end{theorem}

In the parallelogram picture (\cref{fig:vector-ops}), $\vu$, $\vv$ and
$\vu + \vv$ form a triangle, and the theorem says that one side of a triangle
is never longer than the other two together. This is the inequality we met
for numbers in \cref{prop:abs-properties}, now in its original geometric
setting. Equality holds only when $\vu$ and $\vv$ point in the same direction.

\begin{mlnote}
Many systems normalise vectors before comparing them: search engines and
retrieval systems store unit-length embeddings, so that only their directions
matter. Lengths also measure the \emph{size} of things: the length of the
weight vector is used to penalise overly large weights
(\cref{ch:linalg2}), and the length of the gradient decides whether
``gradient clipping'' is needed (\cref{ch:optimization}).
\end{mlnote}

\sectionexercises

\begin{exercise}
Compute $\norm{(3, 4)}$, $\norm{(1, -2, 2)}$ and $\norm{(1, 1, 1, 1)}$.
\end{exercise}
\begin{solution}
$5$; $3$; $\sqrt{4} = 2$.
\end{solution}

\begin{exercise}
Find the unit vector pointing in the same direction as $(6, -8)$.
\end{exercise}
\begin{solution}
$\norm{(6,-8)} = \sqrt{36 + 64} = 10$, so the unit vector is $(0.6, -0.8)$.
\end{solution}

\begin{exercise}
Find the distance between $(1, 2, 3)$ and $(4, 6, 3)$.
\end{exercise}
\begin{solution}
$\norm{(1,2,3) - (4,6,3)} = \norm{(-3, -4, 0)} = 5$.
\end{solution}

\begin{exercise}
Check the triangle inequality for $\vu = (1, 0)$ and $\vv = (0, 1)$. Then find
two non-zero vectors for which it holds with equality.
\end{exercise}
\begin{solution}
$\norm{\vu + \vv} = \norm{(1,1)} = \sqrt2 \approx 1.41 \le 1 + 1 = 2$. Equality
holds for vectors pointing the same way, e.g.\ $\vu = (1, 0)$, $\vv = (2, 0)$:
$\norm{(3, 0)} = 3 = 1 + 2$.
\end{solution}

%----------------------------------------------------------------------
\section{The dot product}\label{sec:dot-product}
%----------------------------------------------------------------------

\begin{motivation}
The dot product is the most important operation in this chapter. It multiplies
two vectors and returns a single number. Algebraically it is a
\emph{weighted sum}, which is exactly what an artificial neuron computes.
Geometrically it measures how much two vectors \emph{point the same way},
which is how machine learning measures similarity.
\end{motivation}

\begin{definition}[Dot product]\label{def:dot-product}
The \term{dot product} (or \term{inner product}) of $\vu, \vv \in \R^n$ is the
number
\[
  \keyeq{\vu \cdot \vv = u_1v_1 + u_2v_2 + \dots + u_nv_n
  = \sum_{i=1}^{n} u_iv_i}
\]
\end{definition}

\begin{example}[Computing a dot product]
$(1, 2, 3)\cdot(4, -5, 6) = 1\cdot4 + 2\cdot(-5) + 3\cdot6 = 4 - 10 + 18 = 12$.
\end{example}

The dot product is commutative ($\vu\cdot\vv = \vv\cdot\vu$), distributes over
addition ($\vu\cdot(\vv + \vw) = \vu\cdot\vv + \vu\cdot\vw$), and lets scalars
move freely ($(c\vu)\cdot\vv = c(\vu\cdot\vv)$). Most importantly, the dot
product of a vector with itself is its squared length:
\[
  \vv \cdot \vv = v_1^2 + \dots + v_n^2 = \norm{\vv}^2 .
\]

\subsection{The geometric meaning}

\begin{theorem}[Geometric formula for the dot product]\label{thm:dot-geometric}
If $\theta$ is the angle between the non-zero vectors $\vu$ and $\vv$, then
\[
  \keyeq{\vu\cdot\vv = \norm{\vu}\,\norm{\vv}\cos\theta}
\]
\end{theorem}

\begin{proof}
The vectors $\vu$, $\vv$ and $\vu - \vv$ form a triangle in which $\theta$ is
the angle between the sides $\vu$ and $\vv$. By the law of cosines
(\cref{thm:law-of-cosines}),
\[
  \norm{\vu - \vv}^2 = \norm{\vu}^2 + \norm{\vv}^2
  - 2\norm{\vu}\,\norm{\vv}\cos\theta .
\]
On the other hand, expanding with the rules of the dot product,
\[
  \norm{\vu - \vv}^2 = (\vu - \vv)\cdot(\vu - \vv)
  = \vu\cdot\vu - 2\,\vu\cdot\vv + \vv\cdot\vv
  = \norm{\vu}^2 - 2\,\vu\cdot\vv + \norm{\vv}^2 .
\]
Comparing the two expressions gives $\vu\cdot\vv = \norm{\vu}\norm{\vv}\cos\theta$.
\end{proof}

Solving for the angle:
\[
  \cos\theta = \frac{\vu\cdot\vv}{\norm{\vu}\,\norm{\vv}} .
\]

\begin{example}[The angle between two vectors]\label{ex:angle-vectors}
Find the angle between $\vu = (2, 1, -1)$ and $\vv = (1, 3, 2)$.

\sol $\vu\cdot\vv = 2 + 3 - 2 = 3$, $\norm{\vu} = \sqrt{6}$ and
$\norm{\vv} = \sqrt{14}$. So $\cos\theta = \frac{3}{\sqrt{84}} \approx 0.327$
and $\theta = \arccos(0.327) \approx 70.9^\circ$.
\end{example}

\begin{intuition}[What the sign of the dot product tells you]
Since lengths are positive, the sign of $\vu\cdot\vv$ is the sign of
$\cos\theta$:
\begin{itemize}
  \item $\vu\cdot\vv > 0$: the angle is less than $90^\circ$; the vectors point
    roughly the same way;
  \item $\vu\cdot\vv = 0$: the angle is exactly $90^\circ$; the vectors are
    perpendicular;
  \item $\vu\cdot\vv < 0$: the angle is more than $90^\circ$; the vectors point
    roughly in opposite directions.
\end{itemize}
\end{intuition}

\begin{definition}[Orthogonal vectors]\label{def:orthogonal}
Two vectors $\vu$ and $\vv$ are \term{orthogonal} (perpendicular), written
$\vu \perp \vv$, if $\vu\cdot\vv = 0$.
\end{definition}

For example, $(1, 2)\cdot(-2, 1) = -2 + 2 = 0$, so these vectors are
orthogonal. This is a remarkably cheap test for a right angle, and it works in
any number of dimensions, where we cannot draw pictures.

\subsection{Projections: the dot product as a shadow}

If $\hat\vv$ is a unit vector, then $\vu\cdot\hat\vv = \norm{\vu}\cos\theta$.
By the right-triangle definition of cosine, this is the length of the
``shadow'' that $\vu$ casts on the line through $\hat\vv$ when light shines
perpendicular to that line (\cref{fig:projection-shadow}). It is called the
\term{scalar projection} of $\vu$ onto the direction $\hat\vv$; it is negative
if the shadow falls backwards. So the dot product with a unit vector answers
the question ``how far does $\vu$ go in the direction $\hat\vv$?''.

\begin{figure}[htbp]
\centering
\begin{tikzpicture}[scale=1.1]
  \coordinate (O) at (0,0);
  \coordinate (U) at (2,3);
  \coordinate (V) at (4,1);
  \coordinate (F) at (2.588,0.647);
  \draw[thick,gray!60] (-0.6,-0.15) -- (5,1.25);
  \draw[vec,PlotRed] (O) -- (V) node[right] {$\vv$};
  \draw[vec,PlotBlue] (O) -- (U) node[above] {$\vu$};
  \draw[dashed,thick] (U) -- (F);
  \draw[line width=3pt,PlotOrange,opacity=0.75] (O) -- (F);
  \pic[draw,angle radius=0.7cm,"$\theta$",angle eccentricity=1.4] {angle=V--O--U};
  \coordinate (F1) at ($(F)!0.22cm!(O)$);
  \coordinate (F2) at ($(F)!0.22cm!(U)$);
  \draw (F1) -- ($(F1)+(F2)-(F)$) -- (F2);
  \draw[decorate,decoration={brace,amplitude=5pt,mirror},thick,PlotOrange]
    ($(O)+(0.08,-0.32)$) -- ($(F)+(0.08,-0.32)$)
    node[midway,below=7pt,font=\small] {$\norm{\vu}\cos\theta$};
\end{tikzpicture}
\caption{The scalar projection of $\vu$ onto the direction of $\vv$: the length
  of the orange ``shadow'' is $\norm{\vu}\cos\theta = \vu\cdot\hat\vv$.}
\label{fig:projection-shadow}
\end{figure}

\begin{theorem}[Cauchy--Schwarz inequality]\label{thm:cauchy-schwarz}
For all $\vu, \vv \in \R^n$, $\ \abs{\vu\cdot\vv} \le \norm{\vu}\,\norm{\vv}$.
\end{theorem}

This follows at once from \cref{thm:dot-geometric}, because
$\abs{\cos\theta} \le 1$. It guarantees that the ``cosine'' computed from the
dot product always lies between $-1$ and $1$, even in a million dimensions.

\subsection{The dot product in machine learning}

\begin{mlnote}[Neurons, similarity and attention]
Dot products are the workhorse of machine learning.
\begin{itemize}
  \item \textbf{Weighted sums.} A linear model, and a single artificial
    neuron before its activation, compute
    \[
      z = w_1x_1 + \dots + w_dx_d + b = \vw\cdot\vx + b .
    \]
    Each weight says how much its feature counts.
  \item \textbf{Cosine similarity.} The \term{cosine similarity}
    $\frac{\vu\cdot\vv}{\norm{\vu}\norm{\vv}}$ compares directions while
    ignoring lengths. Take the made-up three-dimensional ``embeddings''
    \[
      \text{cat} = (0.9, 0.8, 0.1), \quad \text{dog} = (0.8, 0.9, 0.2),
      \quad \text{car} = (0.1, 0.2, 0.9).
    \]
    The similarity of cat and dog is
    $\frac{1.46}{1.208 \times 1.221} \approx 0.99$, while that of cat and car
    is $\frac{0.34}{1.208 \times 0.927} \approx 0.30$.
  \item \textbf{Attention.} In a transformer, each word produces a ``query''
    vector $\vq$ and a ``key'' vector $\vk$, and the relevance of word $j$ to
    word $i$ is measured by the dot product $\vq_i \cdot \vk_j$. Attention is,
    at heart, a large table of dot products.
\end{itemize}
\end{mlnote}

\begin{proposition}[The weight vector is perpendicular to the decision boundary]
\label{prop:normal-vector}
Let $\vw \neq \vzero$. The vector $\vw$ is perpendicular to the decision
boundary $\{\vx : \vw\cdot\vx + b = 0\}$ of a linear classifier.
\end{proposition}

\begin{proof}
Take any two points $\vx$ and $\vx'$ on the boundary. Then
$\vw\cdot\vx + b = 0$ and $\vw\cdot\vx' + b = 0$; subtracting,
$\vw\cdot(\vx - \vx') = 0$. So $\vw$ is orthogonal to every direction
$\vx - \vx'$ that lies within the boundary.
\end{proof}

This explains the picture of \cref{fig:two-classes}: there
$\vw = (1, 1)$, which points from the region of class A straight across the
boundary into the region of class B. The quantity $\vw\cdot\vx + b$ increases
fastest in the direction of $\vw$.

\sectionexercises

\begin{exercise}
Compute $(2, -1, 3)\cdot(4, 0, -2)$ and $(1, 1, 1, 1)\cdot(1, -1, 1, -1)$. Are
either of these pairs orthogonal?
\end{exercise}
\begin{solution}
$8 + 0 - 6 = 2$ (not orthogonal) and $1 - 1 + 1 - 1 = 0$ (orthogonal).
\end{solution}

\begin{exercise}
Find the angle between (a) $(1, 1)$ and $(-1, 0)$; (b) $(1, 2, 2)$ and
$(2, -2, 1)$.
\end{exercise}
\begin{solution}
(a) $\cos\theta = \tfrac{-1}{\sqrt2\cdot1} = -\tfrac{1}{\sqrt2}$, so
$\theta = 135^\circ$. (b) The dot product is $2 - 4 + 2 = 0$, so
$\theta = 90^\circ$.
\end{solution}

\begin{exercise}
For which value of $k$ is $(k, 2)$ orthogonal to $(3, -6)$?
\end{exercise}
\begin{solution}
$3k - 12 = 0$, so $k = 4$.
\end{solution}

\begin{exercise}
Using the made-up embeddings in the box above, compute the cosine similarity of
``dog'' and ``car''. Which pair of the three words is least similar?
\end{exercise}
\begin{solution}
$\text{dog}\cdot\text{car} = 0.08 + 0.18 + 0.18 = 0.44$, and
$\norm{\text{dog}} \approx 1.221$, $\norm{\text{car}} \approx 0.927$, so the
similarity is $\tfrac{0.44}{1.132} \approx 0.39$. The least similar pair is
cat and car ($\approx 0.30$), and the most similar is cat and dog
($\approx 0.99$).
\end{solution}

\begin{exercise}
For the line $3x_1 + 4x_2 - 12 = 0$, give a vector perpendicular to it. Then
find the point of the line closest to the origin, and its distance from the
origin. \emph{Hint:} the closest point lies in the direction of the
perpendicular vector from the origin, so it has the form $t(3, 4)$.
\end{exercise}
\begin{solution}
$\vw = (3, 4)$ is perpendicular to the line (\cref{prop:normal-vector}).
Substituting $(3t, 4t)$ into the equation: $9t + 16t = 12$, so $t = 0.48$ and
the closest point is $(1.44, 1.92)$. Its distance from the origin is
$0.48 \times \norm{(3,4)} = 0.48 \times 5 = 2.4$. (In general the distance
from the origin to $\vw\cdot\vx + b = 0$ is $\abs{b}/\norm{\vw}$.)
\end{solution}

\begin{exercise}
A neuron has weights $\vw = (0.5, -1, 2)$, bias $b = 0.1$ and a sigmoid
activation. Compute its output for the input $\vx = (2, 1, 0.5)$.
\end{exercise}
\begin{solution}
$z = \vw\cdot\vx + b = 1 - 1 + 1 + 0.1 = 1.1$, and the output is
$\sigma(1.1) = \tfrac{1}{1 + e^{-1.1}} \approx \tfrac{1}{1.333} \approx 0.750$.
\end{solution}

%----------------------------------------------------------------------
\section{Matrices}\label{sec:matrices}
%----------------------------------------------------------------------

\begin{definition}[Matrix]\label{def:matrix}
A \term{matrix} is a rectangular array of numbers. A matrix with $m$ rows and
$n$ columns is called an \emph{$m\times n$ matrix} (read ``$m$ by $n$''), and
the set of all such matrices is written $\R^{m\times n}$. The entry in row $i$
and column $j$ of a matrix $\mA$ is written $a_{ij}$ (or $A_{ij}$):
\[
  \mA = \begin{pmatrix}
    a_{11} & a_{12} & \cdots & a_{1n}\\
    a_{21} & a_{22} & \cdots & a_{2n}\\
    \vdots & \vdots & \ddots & \vdots\\
    a_{m1} & a_{m2} & \cdots & a_{mn}
  \end{pmatrix}.
\]
\end{definition}

\begin{notation}
Matrices are written in bold upper-case letters: $\mA$, $\mW$, $\mX$. The
size is always ``rows $\times$ columns'', and the first index of an entry is
always the row. A vector in $\R^n$ is an $n \times 1$ matrix (a single
column). A matrix can be viewed as a list of columns or as a stack of rows;
both views will be useful.
\end{notation}

For example,
\[
  \mA = \begin{pmatrix} 3 & 0 & -1\\ 2 & 5 & 4 \end{pmatrix}
\]
is a $2\times3$ matrix with $a_{13} = -1$ and $a_{22} = 5$. Its second row is
$(2, 5, 4)$ and its third column is $(-1, 4)$.

\begin{mlnote}[Data sets and images are matrices]
A data set of $n$ examples with $d$ features each is stored as an $n\times d$
\term{data matrix} $\mX$: one row per example, one column per feature. For
three houses with features (size, bedrooms),
\[
  \mX = \begin{pmatrix} 100 & 3\\ 80 & 2\\ 120 & 4 \end{pmatrix}.
\]
A grayscale image is a matrix of pixel intensities. A colour image is three
such matrices (red, green, blue) stacked together; arrays with more than two
indices like this are called \term[tensor]{tensors}, which is where the
software library TensorFlow gets its name.
\end{mlnote}

\subsection{Special matrices}

\begin{itemize}
  \item A \emph{square} matrix has as many rows as columns.
  \item The \term{zero matrix} $\mZero$ has all entries equal to $0$.
  \item A \term{diagonal matrix} is square with zeros everywhere off the main
    diagonal (the entries $a_{11}, a_{22}, \dots$), for example
    $\left(\begin{smallmatrix} 2 & 0\\ 0 & -1\end{smallmatrix}\right)$.
  \item The \term{identity matrix} $\mI$ (or $\mI_n$) is the diagonal matrix
    with all diagonal entries equal to $1$.
  \item A square matrix is \term[symmetric matrix]{symmetric} if
    $a_{ij} = a_{ji}$ for all $i, j$, i.e.\ it is a mirror image of itself
    across the main diagonal, like
    $\left(\begin{smallmatrix} 1 & 2\\ 2 & 3\end{smallmatrix}\right)$.
  \item A square matrix is \emph{upper triangular} if all entries below the
    main diagonal are zero.
\end{itemize}

\subsection{Adding and scaling matrices}

Matrices of the same size are added, and multiplied by scalars, entry by
entry, exactly like vectors:
\[
  \begin{pmatrix} 1 & 2\\ 3 & 4\end{pmatrix}
  + 2\begin{pmatrix} 0 & 1\\ 1 & 0\end{pmatrix}
  = \begin{pmatrix} 1 & 4\\ 5 & 4\end{pmatrix}.
\]

\sectionexercises

\begin{exercise}
For $\mA = \left(\begin{smallmatrix} 3 & 0 & -1\\ 2 & 5 & 4\end{smallmatrix}\right)$,
give the size of $\mA$, the entries $a_{21}$ and $a_{13}$, the first row and
the second column.
\end{exercise}
\begin{solution}
$2\times3$; $a_{21} = 2$; $a_{13} = -1$; first row $(3, 0, -1)$; second
column $(0, 5)$.
\end{solution}

\begin{exercise}
Which of these matrices are symmetric?
$\left(\begin{smallmatrix} 1 & 2\\ 2 & 3\end{smallmatrix}\right)$,
$\left(\begin{smallmatrix} 0 & 1\\ -1 & 0\end{smallmatrix}\right)$,
$\left(\begin{smallmatrix} 4 & 0 & 0\\ 0 & 5 & 0\\ 0 & 0 & 6\end{smallmatrix}\right)$.
\end{exercise}
\begin{solution}
The first and the third (every diagonal matrix is symmetric). The second is
not: $a_{12} = 1$ but $a_{21} = -1$.
\end{solution}

\begin{exercise}
Compute $2\mA - \mB$ for
$\mA = \left(\begin{smallmatrix} 1 & 2\\ 3 & 4\end{smallmatrix}\right)$ and
$\mB = \left(\begin{smallmatrix} 0 & 1\\ 1 & 0\end{smallmatrix}\right)$.
\end{exercise}
\begin{solution}
$\left(\begin{smallmatrix} 2 & 4\\ 6 & 8\end{smallmatrix}\right) -
\left(\begin{smallmatrix} 0 & 1\\ 1 & 0\end{smallmatrix}\right) =
\left(\begin{smallmatrix} 2 & 3\\ 5 & 8\end{smallmatrix}\right)$.
\end{solution}

\begin{exercise}
A data set consists of $500$ grayscale images of $28\times28$ pixels. Each
image is flattened into a row of the data matrix. What is the size of the
data matrix?
\end{exercise}
\begin{solution}
Each image has $28 \times 28 = 784$ pixels, so the data matrix is
$500 \times 784$.
\end{solution}

%----------------------------------------------------------------------
\section{Multiplying a matrix by a vector}\label{sec:matrix-vector}
%----------------------------------------------------------------------

\begin{definition}[Matrix--vector product]\label{def:matrix-vector}
Let $\mA$ be an $m\times n$ matrix and $\vx \in \R^n$. The product $\mA\vx$ is
the vector in $\R^m$ whose $i$-th component is the dot product of the $i$-th
row of $\mA$ with $\vx$:
\[
  (\mA\vx)_i = a_{i1}x_1 + a_{i2}x_2 + \dots + a_{in}x_n
  = \sum_{j=1}^{n} a_{ij}x_j .
\]
\end{definition}

The number of columns of $\mA$ must equal the number of components of $\vx$;
the result has as many components as $\mA$ has rows:
\[
  \underbrace{\mA}_{m\times n}\;\underbrace{\vx}_{n\times 1}
  = \underbrace{\mA\vx}_{m\times1}.
\]

\begin{example}[Two ways to compute $\mA\vx$]\label{ex:matrix-vector}
Compute $\mA\vx$ for
$\mA = \left(\begin{smallmatrix} 2 & 1\\ -1 & 3\\ 0 & 4\end{smallmatrix}\right)$
and $\vx = (3, 2)$.

\sol \emph{Row view} (dot products of the rows with $\vx$):
\[
  \begin{pmatrix} 2 & 1\\ -1 & 3\\ 0 & 4\end{pmatrix}
  \begin{pmatrix} 3\\ 2\end{pmatrix}
  = \begin{pmatrix} 2\cdot3 + 1\cdot2\\ -1\cdot3 + 3\cdot2\\ 0\cdot3 + 4\cdot2
    \end{pmatrix}
  = \begin{pmatrix} 8\\ 3\\ 8\end{pmatrix}.
\]
\emph{Column view} (a linear combination of the columns, with the entries of
$\vx$ as coefficients):
\[
  3\begin{pmatrix} 2\\ -1\\ 0\end{pmatrix}
  + 2\begin{pmatrix} 1\\ 3\\ 4\end{pmatrix}
  = \begin{pmatrix} 6\\ -3\\ 0\end{pmatrix}
  + \begin{pmatrix} 2\\ 6\\ 8\end{pmatrix}
  = \begin{pmatrix} 8\\ 3\\ 8\end{pmatrix}.
\]
\end{example}

\begin{intuition}[Rows and columns]
Both views give the same answer, but they tell different stories.
\begin{itemize}
  \item \textbf{Row view:} $\mA\vx$ is a list of $m$ dot products, one for
    each row. If each row is the weight vector of a neuron, $\mA\vx$ computes
    all the neurons of a layer at once.
  \item \textbf{Column view:} $\mA\vx$ is a linear combination of the columns
    of $\mA$, mixed in the amounts given by $\vx$. So the possible outputs
    $\mA\vx$ are exactly the linear combinations of the columns, a fact that
    will be central in \cref{ch:linalg2}.
\end{itemize}
\end{intuition}

Matrix--vector multiplication is \emph{linear}: for all vectors $\vu, \vv$ and
scalars $c$,
\[
  \mA(\vu + \vv) = \mA\vu + \mA\vv, \qquad \mA(c\vu) = c\,\mA\vu ,
\]
which follows directly from the distributive law.

\begin{mlnote}[Predictions for a whole data set, and a neural network layer]
For a linear model $\hat y = \vw\cdot\vx + b$, the predictions for all $n$
examples at once are $\hat\vy = \mX\vw + b$, where $\mX$ is the data matrix
and $b$ is added to every component. With the house data above,
$\vw = (2000, 10\,000)$ and $b = 20\,000$:
\[
  \hat\vy = \begin{pmatrix} 100 & 3\\ 80 & 2\\ 120 & 4 \end{pmatrix}
  \begin{pmatrix} 2000\\ 10\,000 \end{pmatrix} + 20\,000
  = \begin{pmatrix} 250\,000\\ 200\,000\\ 300\,000 \end{pmatrix}.
\]
A \term[layer (neural network)]{layer} of a neural network with weight matrix
$\mW$ (one row per neuron), bias vector $\vb$ and activation $f$ computes
\[
  \keyeq{\vh = f(\mW\vx + \vb)}
\]
where $f$ is applied to each component separately. For example, with
$\mW = \left(\begin{smallmatrix} 1 & -1 & 2\\ 0 & -2 & 1\end{smallmatrix}\right)$,
$\vb = (0.5, -1)$, $\vx = (1, 2, 3)$ and $f = \ReLU$: $\mW\vx + \vb =
(5 + 0.5, -1 - 1) = (5.5, -2)$, and $\vh = (5.5, 0)$.
\end{mlnote}

\sectionexercises

\begin{exercise}
Compute
$\left(\begin{smallmatrix} 1 & 2 & 0\\ 0 & 1 & -1\end{smallmatrix}\right)
\left(\begin{smallmatrix} 1\\ 2\\ 3\end{smallmatrix}\right)$.
\end{exercise}
\begin{solution}
$(1 + 4 + 0,\ 0 + 2 - 3) = (5, -1)$.
\end{solution}

\begin{exercise}
Compute $\left(\begin{smallmatrix} 2 & 1\\ 1 & 3\end{smallmatrix}\right)
\left(\begin{smallmatrix} 1\\ -1\end{smallmatrix}\right)$ using the row view and
again using the column view.
\end{exercise}
\begin{solution}
Rows: $(2 - 1,\ 1 - 3) = (1, -2)$. Columns: $1\cdot(2, 1) + (-1)\cdot(1, 3) =
(2 - 1, 1 - 3) = (1, -2)$.
\end{solution}

\begin{exercise}
Let $\mA$ be a $3\times2$ matrix, $\vx \in \R^3$ and $\vy \in \R^2$. Which of
$\mA\vx$ and $\mA\vy$ is defined, and in which space does it live?
\end{exercise}
\begin{solution}
$\mA\vy$ is defined (the number of columns of $\mA$, $2$, matches the length of
$\vy$) and lies in $\R^3$. $\mA\vx$ is not defined.
\end{solution}

\begin{exercise}
A linear model has weights $\vw = (2, -1)$ and bias $b = 0.5$. Compute its
predictions for the data matrix
$\mX = \left(\begin{smallmatrix} 1 & 2\\ 3 & 0\\ -1 & 1\end{smallmatrix}\right)$.
\end{exercise}
\begin{solution}
$\mX\vw = (2 - 2,\ 6 - 0,\ -2 - 1) = (0, 6, -3)$; adding $b$ gives
$\hat\vy = (0.5, 6.5, -2.5)$.
\end{solution}

%----------------------------------------------------------------------
\section{Matrices as linear transformations}\label{sec:linear-transformations}
%----------------------------------------------------------------------

A matrix $\mA$ of size $m\times n$ turns each vector $\vx\in\R^n$ into a vector
$\mA\vx \in \R^m$, so it defines a function $\vx \mapsto \mA\vx$. Thinking of
matrices as functions that move points around, rather than as tables of
numbers, is the key to understanding them.

\begin{definition}[Linear transformation]\label{def:linear-transformation}
A function $T\colon \R^n \to \R^m$ is a \term{linear transformation} (or
\emph{linear map}) if for all vectors $\vu, \vv$ and scalars $c$,
\[
  T(\vu + \vv) = T(\vu) + T(\vv) \qquad\text{and}\qquad T(c\vu) = c\,T(\vu).
\]
\end{definition}

Every matrix defines a linear transformation, as we just saw. Remarkably, the
converse is also true, and the proof tells us how to find the matrix.

\begin{theorem}[The columns of a matrix are the images of the basis vectors]
\label{thm:columns-images}
For any matrix $\mA$, the $j$-th column of $\mA$ is $\mA\ve_j$. Consequently,
every linear transformation $T\colon\R^n\to\R^m$ is given by a matrix, namely
the matrix whose columns are $T(\ve_1), \dots, T(\ve_n)$.
\end{theorem}

\begin{proof}
By the column view, $\mA\ve_j$ is the combination of columns with coefficient
$1$ for column $j$ and $0$ for all others, which is column $j$ itself. Now let
$T$ be linear and let $\mA$ have columns $T(\ve_1), \dots, T(\ve_n)$. Every
$\vx$ is $x_1\ve_1 + \dots + x_n\ve_n$, so by linearity
$T(\vx) = x_1T(\ve_1) + \dots + x_nT(\ve_n)$, which by the column view is
$\mA\vx$.
\end{proof}

\begin{intuition}[Watch where the basis vectors go]
To understand what a $2\times2$ matrix does, look at its two columns: they show
where $\ve_1 = (1, 0)$ and $\ve_2 = (0, 1)$ end up. Everything else follows,
because every vector is a combination of $\ve_1$ and $\ve_2$ and linearity
preserves combinations. The unit square spanned by $\ve_1$ and $\ve_2$ becomes
the parallelogram spanned by the two columns. Grid lines stay straight,
parallel and evenly spaced, and the origin stays put. \Cref{fig:transformations}
shows a gallery.
\end{intuition}

\begin{figure}[htbp]
\centering
\newcommand{\MLtransformpanel}[5]{%
  \begin{tikzpicture}[scale=0.62]
    \draw[gray!22] (-1.5,-1.5) grid (2.5,2.5);
    \draw[-{Stealth[length=1.5mm]},gray] (-1.5,0) -- (2.7,0);
    \draw[-{Stealth[length=1.5mm]},gray] (0,-1.5) -- (0,2.7);
    \draw[dashed,thick,gray] (0,0) rectangle (1,1);
    \fill[MLdef!20] (0,0) -- (#1,#3) -- ({#1+#2},{#3+#4}) -- (#2,#4) -- cycle;
    \draw[MLdef,thick] (0,0) -- (#1,#3) -- ({#1+#2},{#3+#4}) -- (#2,#4) -- cycle;
    \draw[-{Stealth[length=2mm]},very thick,PlotRed] (0,0) -- (#1,#3);
    \draw[-{Stealth[length=2mm]},very thick,PlotBlue] (0,0) -- (#2,#4);
    \node[font=\footnotesize,align=center,anchor=north] at (0.5,-1.6) {#5};
  \end{tikzpicture}}
\MLtransformpanel{1.5}{0}{0}{0.5}{scaling\\$\left(\begin{smallmatrix}1.5&0\\0&0.5\end{smallmatrix}\right)$}\hfill
\MLtransformpanel{0.866}{-0.5}{0.5}{0.866}{rotation by $30^\circ$\\$\left(\begin{smallmatrix}0.87&-0.5\\0.5&0.87\end{smallmatrix}\right)$}\hfill
\MLtransformpanel{1}{1}{0}{1}{shear\\$\left(\begin{smallmatrix}1&1\\0&1\end{smallmatrix}\right)$}

\medskip
\MLtransformpanel{-1}{0}{0}{1}{reflection\\$\left(\begin{smallmatrix}-1&0\\0&1\end{smallmatrix}\right)$}\hfill
\MLtransformpanel{1}{0}{0}{0.02}{projection onto $x_1$-axis\\$\left(\begin{smallmatrix}1&0\\0&0\end{smallmatrix}\right)$}\hfill
\MLtransformpanel{2}{-0.5}{1}{1.5}{general\\$\left(\begin{smallmatrix}2&-0.5\\1&1.5\end{smallmatrix}\right)$}
\caption{What $2\times2$ matrices do to the unit square (dashed). The red
  arrow is the image of $\ve_1$ (the first column), the blue arrow the image of
  $\ve_2$ (the second column), and the shaded parallelogram is the image of the
  unit square. The projection squashes the square flat onto the $x_1$-axis.}
\label{fig:transformations}
\end{figure}

\begin{example}[Finding the matrix of a transformation]
Find the matrix that rotates the plane anticlockwise by $90^\circ$, and apply
it to $(2, 1)$.

\sol Rotating by $90^\circ$ sends $\ve_1 = (1, 0)$ to $(0, 1)$ and
$\ve_2 = (0, 1)$ to $(-1, 0)$. These are the columns:
\[
  \mR = \begin{pmatrix} 0 & -1\\ 1 & 0\end{pmatrix},
  \qquad
  \mR\begin{pmatrix} 2\\ 1\end{pmatrix} = \begin{pmatrix} -1\\ 2\end{pmatrix},
\]
matching the rotation formula of \cref{thm:rotation}. In general, rotation by
$\alpha$ is given by the \term{rotation matrix}
\[
  \mR_\alpha = \begin{pmatrix} \cos\alpha & -\sin\alpha\\
  \sin\alpha & \cos\alpha\end{pmatrix}.
\]
\end{example}

\begin{remark}[Linear versus affine]
A linear transformation always sends $\vzero$ to $\vzero$, because
$T(\vzero) = T(0\cdot\vzero) = 0\cdot T(\vzero) = \vzero$. So a
\emph{translation} $\vx \mapsto \vx + \vb$ with $\vb \neq \vzero$ is not linear.
Maps of the form $\vx \mapsto \mA\vx + \vb$ are called \term[affine
map]{affine}: a linear transformation followed by a shift. A neural network
layer before its activation function is affine; in machine learning this
distinction is usually blurred and both are called ``linear''.
\end{remark}

\begin{mlnote}[A neural network transforms space]
Each layer of a neural network applies an affine map (which can rotate,
stretch, shear and squash the data cloud) followed by a non-linear activation
(which bends it). Layer after layer, a network can move tangled-up classes of
data into positions where a simple linear boundary separates them. This
geometric picture of ``deforming space until the classes separate'' is one of
the best ways to think about what deep networks do.
\end{mlnote}

\sectionexercises

\begin{exercise}
Find the $2\times2$ matrix that (a) reflects every point in the line $x_2 =
x_1$; (b) doubles every vector; (c) rotates by $180^\circ$.
\end{exercise}
\begin{solution}
(a) $\ve_1 \mapsto (0,1)$ and $\ve_2 \mapsto (1, 0)$, so
$\left(\begin{smallmatrix} 0 & 1\\ 1 & 0\end{smallmatrix}\right)$, which swaps
the two coordinates. (b) $2\mI = \left(\begin{smallmatrix} 2 & 0\\ 0 &
2\end{smallmatrix}\right)$. (c) $\left(\begin{smallmatrix} -1 & 0\\ 0 &
-1\end{smallmatrix}\right)$, i.e.\ $-\mI$; this also equals $\mR_{180^\circ}$.
\end{solution}

\begin{exercise}
Is $T(x_1, x_2) = (x_1 + 1, x_2)$ a linear transformation? Is
$S(x_1, x_2) = (x_2, 3x_1)$? If so, give its matrix.
\end{exercise}
\begin{solution}
$T$ is not linear, since $T(\vzero) = (1, 0) \neq \vzero$ (it is a
translation). $S$ is linear, with matrix $\left(\begin{smallmatrix} 0 & 1\\ 3 &
0\end{smallmatrix}\right)$ (columns $S(\ve_1) = (0, 3)$ and
$S(\ve_2) = (1, 0)$).
\end{solution}

\begin{exercise}
A linear transformation sends $\ve_1$ to $(2, 1)$ and $\ve_2$ to $(-1, 3)$.
Write down its matrix and find the image of $(4, 2)$.
\end{exercise}
\begin{solution}
$\mA = \left(\begin{smallmatrix} 2 & -1\\ 1 & 3\end{smallmatrix}\right)$ and
$\mA(4, 2) = 4(2, 1) + 2(-1, 3) = (6, 10)$.
\end{solution}

%----------------------------------------------------------------------
\section{Matrix multiplication}\label{sec:matrix-multiplication}
%----------------------------------------------------------------------

\subsection{Composing transformations}

What happens if we apply one linear transformation and then another? First
$\vx \mapsto \mB\vx$, then $\mA$ applied to the result: $\vx \mapsto
\mA(\mB\vx)$. A composition of linear maps is linear, so by
\cref{thm:columns-images} it is given by a single matrix. We define the
\emph{product} $\mA\mB$ to be exactly that matrix, so that
\[
  (\mA\mB)\vx = \mA(\mB\vx) \quad\text{for every } \vx .
\]
Its $j$-th column is the image of $\ve_j$, namely $\mA(\mB\ve_j) = \mA\vb_j$,
where $\vb_j$ is the $j$-th column of $\mB$.

\begin{definition}[Matrix product]\label{def:matrix-product}
Let $\mA$ be $m\times n$ and $\mB$ be $n\times p$, with columns
$\vb_1, \dots, \vb_p$. Their \term[matrix multiplication]{product} is the
$m\times p$ matrix
\[
  \mA\mB = \begin{pmatrix} \mA\vb_1 & \mA\vb_2 & \cdots & \mA\vb_p\end{pmatrix},
\]
whose $(i, j)$ entry is the dot product of row $i$ of $\mA$ with column $j$ of
$\mB$:
\[
  (\mA\mB)_{ij} = \sum_{k=1}^{n} a_{ik}b_{kj} .
\]
\end{definition}

\begin{pitfall}[Shapes must match]
$\mA\mB$ is defined only when the number of \emph{columns} of $\mA$ equals the
number of \emph{rows} of $\mB$:
\[
  \underbrace{\mA}_{m\times \boxed{n}}\;\underbrace{\mB}_{\boxed{n}\times p}
  = \underbrace{\mA\mB}_{m\times p}.
\]
The inner sizes must agree and the outer sizes give the size of the result.
Checking shapes like this catches a large fraction of bugs in machine learning
code, where ``shape mismatch'' is the most common error message of all.
\end{pitfall}

\begin{example}[Matrix multiplication is not commutative]\label{ex:not-commutative}
Let $\mA = \left(\begin{smallmatrix} 1 & 2\\ 3 & 4\end{smallmatrix}\right)$ and
$\mB = \left(\begin{smallmatrix} 2 & 0\\ 1 & 3\end{smallmatrix}\right)$. Compute
$\mA\mB$ and $\mB\mA$.

\sol Each entry is ``row of the left matrix $\cdot$ column of the right
matrix'':
\[
  \mA\mB = \begin{pmatrix} 1\cdot2 + 2\cdot1 & 1\cdot0 + 2\cdot3\\
  3\cdot2 + 4\cdot1 & 3\cdot0 + 4\cdot3\end{pmatrix}
  = \begin{pmatrix} 4 & 6\\ 10 & 12\end{pmatrix},
  \qquad
  \mB\mA = \begin{pmatrix} 2 & 4\\ 10 & 14\end{pmatrix}.
\]
So $\mA\mB \neq \mB\mA$.
\end{example}

\begin{intuition}[Why order matters]
$\mA\mB$ means ``first $\mB$, then $\mA$'' (read right to left, like
$f(g(x))$). Doing things in a different order usually gives a different
result: rotating a picture and then stretching it horizontally is not the same
as stretching it horizontally and then rotating it. Putting on socks and then
shoes is not the same as shoes and then socks.
\end{intuition}

\begin{theorem}[Rules of matrix multiplication]\label{thm:matrix-mult-rules}
Whenever the products are defined:
\begin{enumerate}
  \item $(\mA\mB)\mC = \mA(\mB\mC)$ (associative);
  \item $\mA(\mB + \mC) = \mA\mB + \mA\mC$ and $(\mA + \mB)\mC = \mA\mC + \mB\mC$
    (distributive);
  \item $(c\mA)\mB = c(\mA\mB) = \mA(c\mB)$;
  \item $\mA\mI = \mI\mA = \mA$ (the identity does nothing);
  \item but in general $\mA\mB \neq \mB\mA$.
\end{enumerate}
\end{theorem}

Associativity is clear from the composition picture: applying $\mC$, then
$\mB$, then $\mA$ is the same process however we group it. We write $\mA^2$
for $\mA\mA$, $\mA^3$ for $\mA\mA\mA$, and so on.

\begin{pitfall}[Familiar rules that fail for matrices]
Besides commutativity, another rule of ordinary numbers fails: $\mA\mB = \mZero$
does \emph{not} imply that $\mA = \mZero$ or $\mB = \mZero$. For example,
\[
  \begin{pmatrix} 0 & 1\\ 0 & 0\end{pmatrix}
  \begin{pmatrix} 0 & 1\\ 0 & 0\end{pmatrix}
  = \begin{pmatrix} 0 & 0\\ 0 & 0\end{pmatrix}.
\]
Also, $(\mA + \mB)^2 = \mA^2 + \mA\mB + \mB\mA + \mB^2$, which equals
$\mA^2 + 2\mA\mB + \mB^2$ only if $\mA\mB = \mB\mA$.
\end{pitfall}

\begin{mlnote}[Why neural networks need non-linear activations]
Two layers without an activation function in between compute
$\mW_2(\mW_1\vx) = (\mW_2\mW_1)\vx$: a single linear map with matrix
$\mW_2\mW_1$. However many linear layers we stack, the result collapses into
one linear layer, so depth would add nothing. This is the many-input version of
\cref{prop:compose-linear}, and it is why every layer of a neural network ends
with a non-linear activation function.

Matrix products also make computation efficient. To push a whole batch of
inputs (the rows of $\mX$) through a layer at once, compute $\mX\mW^\top$ (the
transpose $\top$ is explained in the next section): one big matrix product
instead of many small ones, which is exactly what GPUs are good at. Multiplying
an $m\times n$ by an $n\times p$ matrix takes $mnp$ multiplications, which is
why the size of these matrices dominates the cost of training large models.
\end{mlnote}

\sectionexercises

\begin{exercise}
Compute $\mA\mB$ and $\mB\mA$ for
$\mA = \left(\begin{smallmatrix} 1 & 0\\ 2 & 1\end{smallmatrix}\right)$ and
$\mB = \left(\begin{smallmatrix} 3 & 1\\ 0 & 2\end{smallmatrix}\right)$.
\end{exercise}
\begin{solution}
$\mA\mB = \left(\begin{smallmatrix} 3 & 1\\ 6 & 4\end{smallmatrix}\right)$ and
$\mB\mA = \left(\begin{smallmatrix} 5 & 1\\ 4 & 2\end{smallmatrix}\right)$.
\end{solution}

\begin{exercise}
Let $\mA$ be $2\times3$, $\mB$ be $3\times4$ and $\mC$ be $4\times2$. Which of
the products $\mA\mB$, $\mB\mC$, $\mC\mA$, $\mA\mC$, $\mA\mB\mC$ are defined, and
what are their sizes?
\end{exercise}
\begin{solution}
$\mA\mB$: $2\times4$. $\mB\mC$: $3\times2$. $\mC\mA$: $4\times3$. $\mA\mC$: not
defined ($3 \neq 4$). $\mA\mB\mC$: $2\times2$.
\end{solution}

\begin{exercise}
Two linear layers have weight matrices
$\mW_1 = \left(\begin{smallmatrix} 1 & 2\\ 0 & 1\end{smallmatrix}\right)$ and
$\mW_2 = \begin{pmatrix} 1 & -1\end{pmatrix}$ (a $1\times2$ matrix). Find the
single matrix $\mW_2\mW_1$ that does the same job, and check that both give the
same output for $\vx = (1, 1)$.
\end{exercise}
\begin{solution}
$\mW_2\mW_1 = \begin{pmatrix} 1\cdot1 + (-1)\cdot0 & 1\cdot2 + (-1)\cdot1
\end{pmatrix} = \begin{pmatrix} 1 & 1\end{pmatrix}$. Step by step:
$\mW_1\vx = (3, 1)$ and $\mW_2(3, 1) = 3 - 1 = 2$. In one go:
$(1, 1)\cdot(1, 1) = 2$. They agree.
\end{solution}

\begin{exercise}
\textbf{(Angle addition formulas.)} Rotating by $\beta$ and then by $\alpha$ is
the same as rotating by $\alpha + \beta$, so $\mR_\alpha\mR_\beta =
\mR_{\alpha + \beta}$. Multiply out $\mR_\alpha\mR_\beta$ and compare entries to
derive the formulas of \cref{thm:angle-addition}.
\end{exercise}
\begin{solution}
\[
  \mR_\alpha\mR_\beta =
  \begin{pmatrix} \cos\alpha\cos\beta - \sin\alpha\sin\beta &
    -\cos\alpha\sin\beta - \sin\alpha\cos\beta\\
    \sin\alpha\cos\beta + \cos\alpha\sin\beta &
    -\sin\alpha\sin\beta + \cos\alpha\cos\beta\end{pmatrix},
\]
while $\mR_{\alpha+\beta}$ has first column
$(\cos(\alpha+\beta), \sin(\alpha+\beta))$. Comparing the first columns gives
$\cos(\alpha+\beta) = \cos\alpha\cos\beta - \sin\alpha\sin\beta$ and
$\sin(\alpha+\beta) = \sin\alpha\cos\beta + \cos\alpha\sin\beta$.
\end{solution}

\begin{exercise}
Show that $\left(\begin{smallmatrix} 1 & 1\\ 0 & 1\end{smallmatrix}\right)^2 =
\left(\begin{smallmatrix} 1 & 2\\ 0 & 1\end{smallmatrix}\right)$, and explain the
result geometrically.
\end{exercise}
\begin{solution}
Row by column: $\left(\begin{smallmatrix} 1\cdot1 + 1\cdot0 & 1\cdot1 +
1\cdot1\\ 0 & 1\end{smallmatrix}\right) = \left(\begin{smallmatrix} 1 & 2\\ 0 &
1\end{smallmatrix}\right)$. Geometrically, the matrix is a horizontal shear
that moves each point sideways by its height; doing it twice moves each point
sideways by twice its height.
\end{solution}

%----------------------------------------------------------------------
\section{The transpose}\label{sec:transpose}
%----------------------------------------------------------------------

\begin{definition}[Transpose]\label{def:transpose}
The \term{transpose} of an $m\times n$ matrix $\mA$ is the $n\times m$ matrix
$\mA^\top$ obtained by turning rows into columns: $(\mA^\top)_{ij} = a_{ji}$.
\end{definition}

For example,
\[
  \begin{pmatrix} 1 & 2 & 3\\ 4 & 5 & 6\end{pmatrix}^{\!\top}
  = \begin{pmatrix} 1 & 4\\ 2 & 5\\ 3 & 6\end{pmatrix}.
\]
A matrix is symmetric exactly when $\mA^\top = \mA$. The transpose of a column
vector is a \emph{row vector}, and this gives a very common way of writing dot
products as matrix products:
\[
  \vu^\top\vv = \begin{pmatrix} u_1 & \cdots & u_n\end{pmatrix}
  \begin{pmatrix} v_1\\ \vdots\\ v_n\end{pmatrix} = \vu\cdot\vv .
\]
(Strictly, $\vu^\top\vv$ is a $1\times1$ matrix, which we identify with a
number.) Machine learning papers write $\vw^\top\vx$ far more often than
$\vw\cdot\vx$.

\begin{theorem}[Rules for transposes]\label{thm:transpose-rules}
$(\mA^\top)^\top = \mA$, $\ (\mA + \mB)^\top = \mA^\top + \mB^\top$,
$\ (c\mA)^\top = c\mA^\top$, and
\[
  (\mA\mB)^\top = \mB^\top\mA^\top .
\]
\end{theorem}

The order reverses in the last rule; you can check it by comparing shapes: if
$\mA$ is $m\times n$ and $\mB$ is $n\times p$, then $(\mA\mB)^\top$ is
$p\times m$, and of the two candidates only $\mB^\top\mA^\top$ ($p\times n$
times $n\times m$) has the right shape.

\begin{proposition}\label{prop:ata-symmetric}
For any matrix $\mA$, the matrix $\mA^\top\mA$ is square and symmetric. Its
$(i,j)$ entry is the dot product of columns $i$ and $j$ of $\mA$.
\end{proposition}

\begin{proof}
$(\mA^\top\mA)^\top = \mA^\top(\mA^\top)^\top = \mA^\top\mA$. The entry formula
follows from ``row $i$ of $\mA^\top$ (which is column $i$ of $\mA$) times column
$j$ of $\mA$''.
\end{proof}

\subsection{The outer product}

Multiplying a column by a row gives a whole matrix. For $\vu \in \R^m$ and
$\vv \in \R^n$, the \term{outer product} $\vu\vv^\top$ is the $m\times n$ matrix
with entries $u_iv_j$:
\[
  \begin{pmatrix} 1\\ 2\end{pmatrix}\begin{pmatrix} 3 & -1 & 0\end{pmatrix}
  = \begin{pmatrix} 3 & -1 & 0\\ 6 & -2 & 0\end{pmatrix}.
\]
Every row is a multiple of $\vv^\top$ and every column is a multiple of $\vu$.
Compare the \emph{inner} product $\vu^\top\vv$ (a number) with the \emph{outer}
product $\vu\vv^\top$ (a matrix).

\begin{mlnote}
The matrices $\mX^\top\mX$ (a $d\times d$ matrix of dot products between
feature columns) and $\mX\mX^\top$ (an $n\times n$ matrix of dot products
between examples) appear everywhere: in the formula for linear regression
(\cref{ch:linalg2}), in covariance matrices (\cref{ch:probability}) and in
kernel methods. Outer products appear in backpropagation: the gradient of a
loss with respect to a weight matrix is an outer product of two vectors
(\cref{ch:dl}).
\end{mlnote}

\sectionexercises

\begin{exercise}
Write down $\mA^\top$ for
$\mA = \left(\begin{smallmatrix} 2 & 0 & -1\\ 1 & 3 & 5\end{smallmatrix}\right)$,
and compute $\mA\mA^\top$.
\end{exercise}
\begin{solution}
$\mA^\top = \left(\begin{smallmatrix} 2 & 1\\ 0 & 3\\ -1 & 5\end{smallmatrix}\right)$
and $\mA\mA^\top = \left(\begin{smallmatrix} 4 + 0 + 1 & 2 + 0 - 5\\ 2 + 0 - 5
& 1 + 9 + 25\end{smallmatrix}\right) = \left(\begin{smallmatrix} 5 & -3\\ -3 &
35\end{smallmatrix}\right)$, which is symmetric, as it should be.
\end{solution}

\begin{exercise}
Using \cref{ex:not-commutative}, verify that $(\mA\mB)^\top = \mB^\top\mA^\top$.
\end{exercise}
\begin{solution}
$(\mA\mB)^\top = \left(\begin{smallmatrix} 4 & 10\\ 6 & 12\end{smallmatrix}\right)$.
And $\mB^\top\mA^\top = \left(\begin{smallmatrix} 2 & 1\\ 0 &
3\end{smallmatrix}\right)\left(\begin{smallmatrix} 1 & 3\\ 2 &
4\end{smallmatrix}\right) = \left(\begin{smallmatrix} 4 & 10\\ 6 &
12\end{smallmatrix}\right)$. They agree.
\end{solution}

\begin{exercise}
For $\vu = (1, 2)$ and $\vv = (3, -1)$, compute $\vu^\top\vv$, $\vu\vv^\top$
and $\vv\vu^\top$.
\end{exercise}
\begin{solution}
$\vu^\top\vv = 3 - 2 = 1$. $\vu\vv^\top = \left(\begin{smallmatrix} 3 & -1\\ 6 &
-2\end{smallmatrix}\right)$ and $\vv\vu^\top = \left(\begin{smallmatrix} 3 & 6\\
-1 & -2\end{smallmatrix}\right) = (\vu\vv^\top)^\top$.
\end{solution}

\begin{exercise}
$\mX$ is an $n\times d$ data matrix. What are the sizes of $\mX^\top\mX$ and
$\mX\mX^\top$? For a data set of a million examples with $100$ features, which
of the two is cheaper to store?
\end{exercise}
\begin{solution}
$\mX^\top\mX$ is $d\times d$ and $\mX\mX^\top$ is $n\times n$. With
$n = 10^6$ and $d = 100$, $\mX^\top\mX$ has $10^4$ entries while $\mX\mX^\top$
has $10^{12}$, so $\mX^\top\mX$ is vastly cheaper.
\end{solution}

%----------------------------------------------------------------------
\section{Inverses and determinants}\label{sec:inverses}
%----------------------------------------------------------------------

\subsection{The inverse of a matrix}

The identity matrix $\mI$ is the transformation that does nothing:
$\mI\vx = \vx$. A transformation can be undone if there is another
transformation that brings every point back to where it started.

\begin{definition}[Inverse matrix]\label{def:inverse-matrix}
A square matrix $\mA$ is \term[invertible matrix]{invertible} if there is a
matrix $\mA^{-1}$, its \term[inverse matrix]{inverse}, such that
\[
  \mA^{-1}\mA = \mA\mA^{-1} = \mI .
\]
A square matrix that has no inverse is called \term[singular
matrix]{singular}.
\end{definition}

Some inverses are obvious from the geometry: the inverse of a rotation by
$\alpha$ is the rotation by $-\alpha$, and the inverse of the scaling
$\left(\begin{smallmatrix} 2 & 0\\ 0 & 0.5\end{smallmatrix}\right)$ is
$\left(\begin{smallmatrix} 0.5 & 0\\ 0 & 2\end{smallmatrix}\right)$. For
$2\times 2$ matrices there is a formula.

\begin{theorem}[Inverse of a $2\times2$ matrix]\label{thm:inverse-2x2}
The matrix $\mA = \left(\begin{smallmatrix} a & b\\ c & d\end{smallmatrix}\right)$
is invertible if and only if $ad - bc \neq 0$, and then
\[
  \mA^{-1} = \frac{1}{ad - bc}\begin{pmatrix} d & -b\\ -c & a\end{pmatrix}.
\]
\end{theorem}

(Swap the diagonal entries, change the signs of the off-diagonal entries,
divide by $ad - bc$.) You can check the formula by multiplying: the product
with $\mA$ is $\frac{1}{ad - bc}\left(\begin{smallmatrix} ad - bc & 0\\ 0 & ad -
bc\end{smallmatrix}\right) = \mI$.

\begin{example}[Inverting a matrix]\label{ex:inverse-2x2}
Find the inverse of $\mA = \left(\begin{smallmatrix} 2 & 1\\ 5 &
3\end{smallmatrix}\right)$.

\sol $ad - bc = 6 - 5 = 1$, so
$\mA^{-1} = \left(\begin{smallmatrix} 3 & -1\\ -5 & 2\end{smallmatrix}\right)$.
\emph{Check:} $\mA\mA^{-1} = \left(\begin{smallmatrix} 6 - 5 & -2 + 2\\ 15 - 15
& -5 + 6\end{smallmatrix}\right) = \mI$. \checkmark
\end{example}

\subsection{The determinant}

The number $ad - bc$ deserves a name and an interpretation.

\begin{definition}[Determinant of a $2\times2$ matrix]\label{def:determinant}
The \term{determinant} of $\mA = \left(\begin{smallmatrix} a & b\\ c &
d\end{smallmatrix}\right)$ is $\det\mA = ad - bc$.
\end{definition}

\begin{intuition}[The determinant measures change of area]
The matrix $\mA$ turns the unit square (area $1$) into the parallelogram
spanned by its columns $(a, c)$ and $(b, d)$, and the area of that
parallelogram is $\abs{ad - bc}$ (\cref{fig:determinant}). Because the
transformation is linear, \emph{every} region gets its area multiplied by the
same factor $\abs{\det\mA}$. The sign records orientation: a negative
determinant means the transformation flips the plane over, like a mirror.
\begin{itemize}
  \item Scaling $\left(\begin{smallmatrix} 2 & 0\\ 0 & 3\end{smallmatrix}\right)$:
    $\det = 6$, areas grow sixfold.
  \item Any rotation: $\det = \cos^2\alpha + \sin^2\alpha = 1$; areas are
    unchanged. A shear also has $\det = 1$.
  \item Reflection $\left(\begin{smallmatrix} -1 & 0\\ 0 & 1\end{smallmatrix}\right)$:
    $\det = -1$; areas unchanged, orientation flipped.
  \item Projection $\left(\begin{smallmatrix} 1 & 0\\ 0 & 0\end{smallmatrix}\right)$:
    $\det = 0$; the plane is squashed onto a line, and areas become $0$.
\end{itemize}
\end{intuition}

\begin{figure}[htbp]
\centering
\begin{tikzpicture}[scale=1.0]
  \draw[gray!22] (-0.5,-0.5) grid (3.5,2.5);
  \draw[-{Stealth[length=2mm]}] (-0.5,0) -- (3.7,0) node[right] {$x_1$};
  \draw[-{Stealth[length=2mm]}] (0,-0.5) -- (0,2.7) node[above] {$x_2$};
  \fill[gray!25] (0,0) rectangle (1,1);
  \node[font=\scriptsize] at (0.5,0.5) {area $1$};
  \fill[MLdef!25,opacity=0.85] (0,0) -- (2,0.5) -- (3,2) -- (1,1.5) -- cycle;
  \draw[MLdef,thick] (0,0) -- (2,0.5) -- (3,2) -- (1,1.5) -- cycle;
  \draw[-{Stealth[length=2mm]},very thick,PlotRed] (0,0) -- (2,0.5) node[right,font=\small] {$(2, 0.5)$};
  \draw[-{Stealth[length=2mm]},very thick,PlotBlue] (0,0) -- (1,1.5) node[above left,font=\small] {$(1, 1.5)$};
  \node[font=\small] at (1.75,1.0) {area $2.5$};
\end{tikzpicture}
\caption{The matrix $\left(\begin{smallmatrix} 2 & 1\\ 0.5 &
  1.5\end{smallmatrix}\right)$ maps the unit square to the parallelogram
  spanned by its columns. Its determinant $2\times1.5 - 1\times0.5 = 2.5$ is
  the area of the parallelogram.}
\label{fig:determinant}
\end{figure}

This picture explains when a matrix is invertible. If $\det\mA = 0$, the
transformation squashes the plane onto a line (or a point). Many different
points land on the same image, so there is no way to undo it: $\mA$ is
singular. For example, the columns of
$\left(\begin{smallmatrix} 1 & 2\\ 2 & 4\end{smallmatrix}\right)$ are parallel
($(2, 4) = 2(1, 2)$), its determinant is $4 - 4 = 0$, and it squashes the
whole plane onto the line through $(1, 2)$.

\subsection{Determinants of larger matrices}

Determinants exist for square matrices of every size, with the same meaning:
for a $3\times3$ matrix, $\abs{\det\mA}$ is the factor by which volumes are
multiplied. A $3 \times 3$ determinant can be computed by \emph{expansion
along the first row}:
\[
  \det\begin{pmatrix} a & b & c\\ d & e & f\\ g & h & i\end{pmatrix}
  = a\det\begin{pmatrix} e & f\\ h & i\end{pmatrix}
  - b\det\begin{pmatrix} d & f\\ g & i\end{pmatrix}
  + c\det\begin{pmatrix} d & e\\ g & h\end{pmatrix}.
\]
Each entry of the first row multiplies the determinant of the $2\times2$
matrix left after deleting that entry's row and column, with alternating signs
$+, -, +$. For example,
\[
  \det\begin{pmatrix} 2 & 0 & 1\\ 1 & 3 & 2\\ 1 & 1 & 2\end{pmatrix}
  = 2(6 - 2) - 0(2 - 2) + 1(1 - 3) = 8 - 0 - 2 = 6 .
\]
For triangular matrices the determinant is simply the product of the diagonal
entries. In practice, computers calculate determinants by elimination
(\cref{sec:linear-systems}), never by expansion, which becomes impossibly slow
for large matrices.

\begin{theorem}[Properties of determinants and inverses]\label{thm:det-properties}
For square matrices $\mA$ and $\mB$ of the same size:
\begin{enumerate}
  \item $\mA$ is invertible if and only if $\det\mA \neq 0$;
  \item $\det(\mA\mB) = \det(\mA)\det(\mB)$ and $\det(\mA^\top) = \det\mA$;
  \item if $\mA$ is invertible, $\det(\mA^{-1}) = 1/\det\mA$;
  \item if $\mA$ and $\mB$ are invertible, $(\mA\mB)^{-1} = \mB^{-1}\mA^{-1}$
    and $(\mA^\top)^{-1} = (\mA^{-1})^\top$.
\end{enumerate}
\end{theorem}

The rule $\det(\mA\mB) = \det\mA\det\mB$ is the area picture again: scaling
areas by $\det\mB$ and then by $\det\mA$ scales them by the product. And to
undo ``first $\mB$, then $\mA$'' you must first undo $\mA$ and then undo
$\mB$, hence the reversed order in $(\mA\mB)^{-1} = \mB^{-1}\mA^{-1}$.

\begin{pitfall}[There is no matrix division]
We never write $\frac{\mA}{\mB}$, because $\mA\mB^{-1}$ and $\mB^{-1}\mA$ are
generally different. Only square matrices can have inverses, and many square
matrices (those with determinant $0$) do not. Finally, even when an inverse
exists, computing it explicitly is slow and can be numerically inaccurate; to
solve $\mA\vx = \vb$, software uses elimination instead
(\cref{sec:linear-systems}).
\end{pitfall}

\begin{mlnote}
Determinants appear in the formula for the multivariate Gaussian distribution
(\cref{ch:probability}), where $\det$ of a covariance matrix measures the
volume of a data cloud. They are also central to \emph{normalising flows}, a
family of generative models that transform simple random noise into complex
data with invertible maps and must keep track of how volumes change.
\end{mlnote}

\sectionexercises

\begin{exercise}
Find the determinant and the inverse of
$\left(\begin{smallmatrix} 4 & 7\\ 2 & 6\end{smallmatrix}\right)$, and check
your answer.
\end{exercise}
\begin{solution}
$\det = 24 - 14 = 10$, so the inverse is
$\tfrac{1}{10}\left(\begin{smallmatrix} 6 & -7\\ -2 & 4\end{smallmatrix}\right)
= \left(\begin{smallmatrix} 0.6 & -0.7\\ -0.2 & 0.4\end{smallmatrix}\right)$.
Check: $\left(\begin{smallmatrix} 4 & 7\\ 2 & 6\end{smallmatrix}\right)
\left(\begin{smallmatrix} 0.6 & -0.7\\ -0.2 & 0.4\end{smallmatrix}\right) =
\left(\begin{smallmatrix} 2.4 - 1.4 & -2.8 + 2.8\\ 1.2 - 1.2 & -1.4 +
2.4\end{smallmatrix}\right) = \mI$.
\end{solution}

\begin{exercise}
For which value of $k$ is $\left(\begin{smallmatrix} k & 2\\ 3 &
6\end{smallmatrix}\right)$ singular? What does it do to the plane for that
value?
\end{exercise}
\begin{solution}
$\det = 6k - 6 = 0$ when $k = 1$. Then both columns, $(1, 3)$ and $(2, 6)$, lie
on the line through $(1, 3)$, and the matrix squashes the whole plane onto that
line.
\end{solution}

\begin{exercise}
Compute $\det\left(\begin{smallmatrix} 1 & 2 & 3\\ 0 & 4 & 5\\ 0 & 0 &
6\end{smallmatrix}\right)$ by expansion along the first row, and compare with
the product of the diagonal entries.
\end{exercise}
\begin{solution}
$1\cdot(24 - 0) - 2\cdot(0 - 0) + 3\cdot(0 - 0) = 24$, which equals
$1\cdot4\cdot6$, as expected for a triangular matrix.
\end{solution}

\begin{exercise}
Let $\mA$ and $\mB$ be $2\times2$ matrices with $\det\mA = 3$ and
$\det\mB = -2$. Find $\det(\mA\mB)$, $\det(\mA^{-1})$ and $\det(2\mA)$.
\end{exercise}
\begin{solution}
$\det(\mA\mB) = 3 \times (-2) = -6$. $\det(\mA^{-1}) = \tfrac13$. For $2\mA$:
the extra factor $2$ doubles lengths in both directions, which multiplies
areas by $4$, so $\det(2\mA) = 4\times3 = 12$. (In general, for an $n\times n$
matrix, $\det(c\mA) = c^n\det\mA$.)
\end{solution}

%----------------------------------------------------------------------
\section{Systems of linear equations}\label{sec:linear-systems}
%----------------------------------------------------------------------

\subsection{From equations to matrices}

In \cref{sec:equations} we solved two equations in two unknowns. With matrices,
any number of linear equations in any number of unknowns can be written as a
single equation. The system
\[
  \begin{cases}
    \phantom{2}x_1 + 2x_2 + \phantom{3}x_3 = 8\\
    2x_1 - \phantom{2}x_2 + 3x_3 = 9\\
    3x_1 + \phantom{2}x_2 - \phantom{3}x_3 = 2
  \end{cases}
  \qquad\text{is}\qquad
  \underbrace{\begin{pmatrix} 1 & 2 & 1\\ 2 & -1 & 3\\ 3 & 1 & -1\end{pmatrix}}_{\mA}
  \underbrace{\begin{pmatrix} x_1\\ x_2\\ x_3\end{pmatrix}}_{\vx}
  = \underbrace{\begin{pmatrix} 8\\ 9\\ 2\end{pmatrix}}_{\vb},
\]
because each row of $\mA\vx$ is the left side of one equation. The matrix
$\mA$ holds the coefficients, $\vx$ the unknowns and $\vb$ the right-hand
sides.

\subsection{Two pictures of a linear system}

\textbf{Row picture.} With two unknowns, each equation describes a line in the
plane, and a solution is a point on all the lines at once. Two lines can meet
in exactly one point, be parallel and never meet, or coincide
(\cref{fig:systems}). So a linear system has either exactly one solution, no
solution, or infinitely many solutions; it can never have exactly two. The
same holds in any dimension, where each equation describes a plane or a
hyperplane.

\textbf{Column picture.} By the column view, $\mA\vx = \vb$ asks: \emph{which
linear combination of the columns of $\mA$ produces $\vb$?} The unknowns are
the amounts of each column in the recipe.

\begin{figure}[htbp]
\centering
\begin{tikzpicture}
\begin{groupplot}[group style={group size=3 by 1,horizontal sep=0.9cm},
  mlplot,width=5.2cm,height=5cm,xmin=-0.5,xmax=4,ymin=-0.5,ymax=4,
  xtick={0,...,3},ytick={0,...,3},title style={font=\small,align=center}]
  \nextgroupplot[title={one solution\\$x+y=3$, $x-y=1$}]
    \addplot[PlotBlue,domain=-0.5:4]{3-x};
    \addplot[PlotRed,domain=-0.5:4]{x-1};
    \addplot[only marks,mark=*,mark size=2pt] coordinates {(2,1)};
  \nextgroupplot[title={no solution\\$x+y=3$, $x+y=1$}]
    \addplot[PlotBlue,domain=-0.5:4]{3-x};
    \addplot[PlotRed,domain=-0.5:4]{1-x};
  \nextgroupplot[title={infinitely many\\$x+y=3$, $2x+2y=6$}]
    \addplot[PlotBlue,line width=3pt,opacity=0.4,domain=-0.5:4]{3-x};
    \addplot[PlotRed,dashed,domain=-0.5:4]{3-x};
\end{groupplot}
\end{tikzpicture}
\caption{The row picture of three systems of two equations in two unknowns:
  lines that cross once, parallel lines, and two equations describing the same
  line.}
\label{fig:systems}
\end{figure}

\subsection{Gaussian elimination}

The standard method for solving linear systems, by hand or by computer, is
\term{Gaussian elimination}. It is the elimination method of
\cref{sec:equations}, organised systematically. We work with the
\term{augmented matrix} $[\mA \mid \vb]$, which records the system without the
variable names, and apply \term[elementary row operation]{elementary row
operations}, none of which changes the solutions:
\begin{enumerate}[label=(\roman*)]
  \item swap two rows;
  \item multiply a row by a non-zero number;
  \item add a multiple of one row to another row.
\end{enumerate}
The goal is a \emph{triangular} (``row echelon'') form in which each equation
has fewer unknowns than the one above it; the last equation is then solved
first and the others by \emph{back substitution}.

\begin{example}[Gaussian elimination]\label{ex:gaussian-elimination}
Solve the $3\times3$ system above.

\sol Start with the augmented matrix and eliminate $x_1$ from rows 2 and 3:
\[
  \left(\begin{array}{rrr|r} 1 & 2 & 1 & 8\\ 2 & -1 & 3 & 9\\ 3 & 1 & -1 & 2
  \end{array}\right)
  \xrightarrow[\;R_3 - 3R_1\;]{\;R_2 - 2R_1\;}
  \left(\begin{array}{rrr|r} 1 & 2 & 1 & 8\\ 0 & -5 & 1 & -7\\ 0 & -5 & -4 & -22
  \end{array}\right).
\]
Now eliminate $x_2$ from row 3:
\[
  \xrightarrow{\;R_3 - R_2\;}
  \left(\begin{array}{rrr|r} 1 & 2 & 1 & 8\\ 0 & -5 & 1 & -7\\ 0 & 0 & -5 & -15
  \end{array}\right).
\]
Back substitution: the last row says $-5x_3 = -15$, so $x_3 = 3$. The second
row says $-5x_2 + x_3 = -7$, so $-5x_2 = -10$ and $x_2 = 2$. The first row says
$x_1 + 2x_2 + x_3 = 8$, so $x_1 = 8 - 4 - 3 = 1$. The solution is
$\vx = (1, 2, 3)$.

\emph{Check} in the original equations: $1 + 4 + 3 = 8$, $2 - 2 + 9 = 9$,
$3 + 2 - 3 = 2$. \checkmark
\end{example}

If elimination produces a row $(0\ 0\ \cdots\ 0 \mid c)$ with $c \neq 0$, that
row says $0 = c$, and the system has no solution. If it produces a row of
zeros $(0\ \cdots\ 0 \mid 0)$, one equation was redundant, and there may be
infinitely many solutions.

\subsection{Solutions, inverses and determinants}

For a square system $\mA\vx = \vb$, the three cases are decided by the
determinant:

\begin{theorem}[Square systems]\label{thm:square-systems}
If $\mA$ is square and $\det\mA \neq 0$, then $\mA\vx = \vb$ has exactly one
solution for every $\vb$, namely $\vx = \mA^{-1}\vb$. If $\det\mA = 0$, the
system has either no solution or infinitely many, depending on $\vb$.
\end{theorem}

Indeed, multiplying $\mA\vx = \vb$ on the left by $\mA^{-1}$ gives
$\mA^{-1}\mA\vx = \mA^{-1}\vb$, i.e.\ $\vx = \mA^{-1}\vb$. This is a useful
formula for thinking, even though (as noted above) software finds $\vx$ by
elimination rather than by computing $\mA^{-1}$.

\begin{mlnote}[More equations than unknowns]
In machine learning, systems are rarely square. Fitting a linear model
$\hat\vy = \mX\vw$ to $n$ data points with $d$ features gives $n$ equations
$\mX\vw = \vy$ in $d$ unknowns, and usually $n$ is much larger than $d$. Such
an \emph{overdetermined} system typically has no exact solution, because no
line passes exactly through a noisy cloud of points. Instead we look for the
$\vw$ that makes $\mX\vw$ as close as possible to $\vy$. This
\emph{least-squares} solution satisfies the square system
\[
  \mX^\top\mX\,\vw = \mX^\top\vy ,
\]
called the \emph{normal equations}, which we derive geometrically in
\cref{ch:linalg2} and with calculus in \cref{ch:optimization}.
\end{mlnote}

\begin{remark}[Linear algebra in code]
In Python, the NumPy library provides all the operations of this chapter. Note
that \texttt{@} is matrix multiplication, while \texttt{*} multiplies entry by
entry.
\par\noindent\begin{minipage}{\linewidth}
\begin{lstlisting}
import numpy as np

A = np.array([[2, 1],
              [-1, 3],
              [0, 4]])         # a 3 x 2 matrix
x = np.array([3, 2])           # a vector in R^2
print(A @ x)                   # [8 3 8]
print(A.T)                     # the 2 x 3 transpose

B = np.array([[1, 2],
              [3, 4]])
print(np.linalg.det(B))        # -2.0 (up to rounding)
print(np.linalg.solve(B, np.array([5, 6])))   # [-4.   4.5]
\end{lstlisting}
\end{minipage}
\end{remark}

\sectionexercises

\begin{exercise}
Write the system $2x - y = 1$, $x + 3y = 11$ in the form $\mA\vx = \vb$ and
solve it using the inverse of $\mA$.
\end{exercise}
\begin{solution}
$\mA = \left(\begin{smallmatrix} 2 & -1\\ 1 & 3\end{smallmatrix}\right)$,
$\vb = (1, 11)$. $\det\mA = 6 + 1 = 7$, so
$\mA^{-1} = \tfrac17\left(\begin{smallmatrix} 3 & 1\\ -1 &
2\end{smallmatrix}\right)$ and $\vx = \tfrac17(3 + 11, -1 + 22) = \tfrac17(14,
21) = (2, 3)$. Check: $4 - 3 = 1$ and $2 + 9 = 11$.
\end{solution}

\begin{exercise}
Solve by Gaussian elimination:
$x_1 + x_2 + x_3 = 6$, $\ 2x_1 + x_2 - x_3 = 1$, $\ x_1 - x_2 + 2x_3 = 5$.
\end{exercise}
\begin{solution}
$R_2 - 2R_1$ gives $(0, -1, -3 \mid -11)$; $R_3 - R_1$ gives
$(0, -2, 1 \mid -1)$. Then $R_3 - 2R_2$ gives $(0, 0, 7 \mid 21)$, so
$x_3 = 3$. From $-x_2 - 3x_3 = -11$: $x_2 = 2$. From the first equation:
$x_1 = 6 - 2 - 3 = 1$. Solution $(1, 2, 3)$.
\end{solution}

\begin{exercise}
How many solutions does each system have? (a) $x + 2y = 4$, $2x + 4y = 8$;
(b) $x + 2y = 4$, $2x + 4y = 5$; (c) $x + 2y = 4$, $x - y = 1$.
\end{exercise}
\begin{solution}
(a) Infinitely many: the second equation is twice the first (the same line).
(b) None: doubling the first gives $2x + 4y = 8 \neq 5$ (parallel lines).
(c) Exactly one: subtracting gives $3y = 3$, so $y = 1$ and $x = 2$.
\end{solution}

\begin{exercise}
We try to fit a line $y = w_1x + w_0$ exactly through the three points
$(1, 2)$, $(2, 3)$ and $(3, 5)$. Write the three equations in matrix form and
show that they have no solution.
\end{exercise}
\begin{solution}
$\left(\begin{smallmatrix} 1 & 1\\ 2 & 1\\ 3 & 1\end{smallmatrix}\right)
\left(\begin{smallmatrix} w_1\\ w_0\end{smallmatrix}\right) =
\left(\begin{smallmatrix} 2\\ 3\\ 5\end{smallmatrix}\right)$. The first two
equations, $w_1 + w_0 = 2$ and $2w_1 + w_0 = 3$, force $w_1 = 1$ and $w_0 = 1$;
but then the third gives $3 + 1 = 4 \neq 5$. The points do not lie on a line,
so the system has no solution; \cref{ch:linalg2} finds the best approximate
one.
\end{solution}

%----------------------------------------------------------------------
\section{Chapter summary}\label{sec:ch4-summary}
%----------------------------------------------------------------------

\begin{keytakeaways}
\begin{itemize}
  \item A vector is both an arrow and a list of numbers. Vectors are added and
    scaled component by component; linear combinations
    $c_1\vv_1 + \dots + c_k\vv_k$ are the central construction.
  \item $\norm{\vv} = \sqrt{\vv\cdot\vv}$ is the length; $\vv/\norm{\vv}$ is a unit
    vector; $\norm{\vu - \vv}$ is a distance.
  \item The dot product $\vu\cdot\vv = \vu^\top\vv = \sum_i u_iv_i =
    \norm{\vu}\norm{\vv}\cos\theta$ is a weighted sum and a measure of
    alignment; $\vu\cdot\vv = 0$ means orthogonal. Cosine similarity and
    neurons are dot products.
  \item An $m\times n$ matrix maps $\R^n$ to $\R^m$. $\mA\vx$ is a list of dot
    products (rows) and a combination of columns. The columns are the images
    of the standard basis vectors.
  \item $\mA\mB$ means ``first $\mB$, then $\mA$''; entries are row $\cdot$
    column; shapes must match; $\mA\mB \neq \mB\mA$ in general. Stacked linear
    layers collapse into one.
  \item $(\mA\mB)^\top = \mB^\top\mA^\top$; $\mA^\top\mA$ is symmetric.
  \item $\det\mA$ is the factor by which areas (volumes) change;
    $\mA^{-1}$ exists exactly when $\det\mA \neq 0$.
  \item $\mA\vx = \vb$ has one, none or infinitely many solutions, found by
    Gaussian elimination. Machine learning systems are usually overdetermined
    and are solved by least squares.
\end{itemize}
\end{keytakeaways}
